\chapter{Numerical Methods and Convergence}
\label{app:numerics}

This appendix states how the linear algebra and the time integration in this
thesis are actually performed, and collects the convergence and conditioning
evidence that has a producing script and a stamped artifact. Results quoted in
the body that lack one are identified as such here rather than repeated.

\section{How the systems are solved}

\paragraph{Steady states.} Every quasi-steady and collisional--radiative
equilibrium population in this thesis is obtained by a direct dense solve of
$\Lmat_{EE}\,\nvec_E = -\Svec_E$ or
$\Lmat_{EE}\,\nvec_E = -\Lmat_{E,g}\,\ngs$, using LAPACK's general solver
through \texttt{numpy.linalg.solve}. No iterative method and no preconditioner
is used, and none is needed: the excited-state block is $42 \times 42$. The two
feed channels are solved separately and their sum checked against the full solve
rather than assumed (Section~\ref{sec:superposition}).

\paragraph{Timescales.} The relaxation spectrum is the eigenvalue spectrum of
$\Lmat$, computed with \texttt{numpy.linalg.eigvals}, sorted descending on the
real part. $\tauslow$ and $\taurel$ are the reciprocals of the two least
negative eigenvalues, and $\Msep = \tauslow/\taurel$.

The selection of ``negative'' eigenvalues is a place where this thesis has
already been bitten once. An earlier version of the filter discarded every
eigenvalue with $\lambda > -1\,\mathrm{s^{-1}}$, on the reasoning that such
values were numerical noise near zero. They are not: at the cold corner of the
grid the true slowest mode has $|\lambda| \approx \SI{0.0149}{\per\second}$,
corresponding to $\tauslow = \SI{67.2}{\second}$, and the filter discarded it and
returned the second mode in its place, wrong by nine orders of magnitude. The
filter is now $\lambda < 0$, which admits every decaying mode. Chapter~\ref{ch:atoms}
reports the error and its effect; the point here is the method, which is that no
magnitude threshold is applied to an eigenvalue whose magnitude is the quantity
being measured.

\paragraph{Time integration.} Transient solutions are obtained two ways, and the
two are compared rather than trusted separately. The spectral route
diagonalises $\Lmat$ and propagates the deviation from the post-step steady
state analytically. The direct route integrates
$\dot{\nvec} = \Lmat\nvec + \Svec\nion$ with the implicit Radau method at
relative tolerance $10^{-10}$ and absolute tolerance
$10^{-12}\max|\nvec_{\rm old}|$, chosen because the system is stiff: the
spectrum spans between eight and fifteen orders of magnitude in $|\lambda|$
depending on the grid point.

\section{Conditioning}

Non-normality matters here more than conditioning does. The operator's
numerical abscissa $\mu(\Lmat)$, which bounds the initial growth rate of
$\|\nvec\|$, is positive everywhere on the grid while the spectral abscissa is
negative everywhere: the solution can grow before it decays, even though every
mode decays. Measured over all $400$ grid points
(\texttt{verify\_operator\_conditioning.py},
\texttt{validation/operator\_conditioning/}):

\begin{table}[htbp]
  \centering
  \caption[Conditioning and non-normality over the grid]{Conditioning of the
    excited-state block and the two abscissae, over all $400$ grid points.}
  \label{tab:conditioning}
  \begin{tabular}{@{}lrrr@{}}
    \toprule
    quantity & minimum & maximum & at the benchmark \\
    \midrule
    $\kappa(\Lmat_{FF})$              & $1.48\times10^{3}$  & $1.74\times10^{5}$  & $4.20\times10^{4}$ \\
    $\mu(\Lmat)$ \si{\per\second}     & $3.40\times10^{8}$  & $1.51\times10^{12}$ & $1.28\times10^{11}$ \\
    spectral abscissa \si{\per\second} & $-1.33\times10^{7}$ & $-1.49\times10^{-2}$ & $-4.40\times10^{4}$ \\
    \bottomrule
  \end{tabular}
\end{table}

\noindent A condition number of $1.7\times10^{5}$ at its worst leaves roughly
eleven significant figures in double precision, against results quoted here to
three or four. Conditioning is not a threat to any number in this thesis. The
positive numerical abscissa is not a numerical problem either. It is a
statement about transient growth of the population vector in the Euclidean
norm of the coordinates used, and would change under a rescaling of those
coordinates; Chapter~\ref{ch:theory} uses it in that sense, and the direct
trajectory tests of Chapter~\ref{ch:validation} are the coordinate-free
evidence.

\section{Convergence tests with a stamped artifact}

\begin{table}[htbp]
  \centering
  \caption[Reproducible numerical checks]{Numerical checks that have a
    producing script and a stamped artifact under \texttt{validation/}.}
  \label{tab:numerical_checks}
  \small
  \begin{tabular}{@{}p{5.2cm}p{3.1cm}p{5.0cm}@{}}
    \toprule
    check & result & artifact \\
    \midrule
    Two-channel superposition: $\nzero + \none$ against the full solve, $784$ pairs
      & $3.08\times10^{-14}$
      & \texttt{divertor\_map/} \\
    Column-sum closure: $\sum_i \Lmat_{ij} = -\nel K_{\rm ion}(j)$, all $43$ states, at two operators
      & $1.75\times10^{-12}$
      & \texttt{fujimoto\_bundle/} \\
    Integrator tolerance at two grid points, $\mathrm{rtol}=10^{-9}$ against $10^{-12}$
      & $1.20\times10^{-8}$
      & \texttt{plateau\_slowmode/} \\
    Direct integration against matrix exponential, $20$ log-spaced times
      & $1.14\times10^{-9}$ at the benchmark, worst $3.24\times10^{-8}$ in $\|\Delta\nvec\|/\|\nvec\|$
      & \texttt{ramp\_plateau/} \\
    Statistical-$\ell$ projector identity $\mathsf{P}\mathsf{R} = \mathsf{I}$
      & $0$ exactly
      & \texttt{fujimoto\_bundle/} \\
    \bottomrule
  \end{tabular}
\end{table}

\noindent Each of these is a check on the numerics, not on the physics. They
establish that the answers reported are the answers the equations give; whether
the equations are the right ones is the subject of Chapter~\ref{ch:validation}
and Chapter~\ref{ch:limits}.

\section{State-space convergence, and what is not established}

Truncation of the level ladder was tested downward, by rebuilding $\Lmat$ at
successively lower $n_{\max}$ with the loss channels of the removed states
returned to the diagonal of the survivors, and the result is reported in
Section~\ref{sec:convergence}. Two limitations of that result belong here rather
than in the body.

First, it is a convergence test of one quantity, $\ffrac{3}-\ffrac{4}$, and not
of the state space in general. Quantities that depend on the Rydberg population
directly, among them the effective ionisation coefficient, are not covered by
it.

Second, the numbers reported in Section~\ref{sec:convergence} were recorded in
this project's working notes and no script in the repository regenerates them.
They are therefore not reproducible in the sense the rest of this appendix uses,
and they are not repeated here. Closing that gap requires writing the truncation
scan as a stamped script; the methodological trap it must avoid is recorded in
Section~\ref{sec:convergence}, since an earlier attempt that dropped states
without returning their loss channels to the diagonal produced a spurious
\SIrange{22}{37}{\percent} jump that looked like a finding.

Convergence \emph{upward}, at $n_{\max} = 18$ or $20$, is not attempted and
cannot be, because the atomic data stop at $n = 15$. That is the second of the
five open questions of Section~\ref{sec:scope_open}.
